\documentclass[11pt,a4paper]{article}
\usepackage[utf8]{inputenc}
\usepackage[T1]{fontenc}
\usepackage[spanish]{babel}
\usepackage{amsmath,amssymb,mathtools}
\usepackage{geometry}
\usepackage{lmodern}
\usepackage{microtype}
\geometry{margin=2.4cm}
\setlength{\parindent}{0pt}
\setlength{\parskip}{0.45em}
\newcommand{\R}{\mathbb{R}}
\newcommand{\sen}{\operatorname{sen}}
\newcommand{\Dom}{\operatorname{Dom}}
\newcommand{\Ima}{\operatorname{Im}}
\newcommand{\Int}{\operatorname{int}}
\newcommand{\Ext}{\operatorname{ext}}
\newcommand{\Fr}{\operatorname{fr}}
\newcommand{\Adh}{\operatorname{adh}}
\newcommand{\Ac}{\operatorname{ac}}
\begin{document}

\begin{center}
{\LARGE\bfseries Análisis Matemático II}\\[0.4em]
{\Large Tema 1: Conceptos fundamentales de funciones de varias variables}\\[0.4em]
{\large Soluciones}
\end{center}

\section{Ejercicio 1}
Sean
\[
u=(-1,1,0,-2,0),\qquad v=(3,-1,2,2,3).
\]
\subsection{Solución}
\[
u\cdot v=-3-1-4=-8.
\]
\[
\boxed{u\cdot v=-8}
\]
\[
\|u\|=\sqrt{1+1+4}=\sqrt6,
\qquad
\|v\|=\sqrt{9+1+4+4+9}=3\sqrt3.
\]
\[
\boxed{\|u\|=\sqrt6,\qquad \|v\|=3\sqrt3}
\]
\[
d(u,v)=\|u-v\|
=\|(-4,2,-2,-4,-3)\|
=\sqrt{49}=7.
\]
\[
\boxed{d(u,v)=7}
\]
\[
u+v=(2,0,2,0,3),
\qquad
\|u+v\|=\sqrt{17}\leq \sqrt6+3\sqrt3.
\]
\[
\boxed{\|u+v\|\leq\|u\|+\|v\|}
\]
\[
|u\cdot v|=8\leq \sqrt6\,3\sqrt3=9\sqrt2.
\]
\[
\boxed{|u\cdot v|\leq\|u\|\|v\|}
\]

\section{Ejercicio 2}
\subsection{Solución}
Si \(y=0\), la desigualdad es inmediata. Si \(y\neq0\),
\[
\left\|x-\lambda y\right\|^2\geq0.
\]
Tomando
\[
\lambda=\frac{x\cdot y}{\|y\|^2},
\]
se obtiene
\[
0\leq
\|x\|^2-\frac{(x\cdot y)^2}{\|y\|^2}.
\]
Por tanto,
\[
(x\cdot y)^2\leq \|x\|^2\|y\|^2
\]
y
\[
\boxed{|x\cdot y|\leq\|x\|\|y\|}.
\]

\section{Ejercicio 3}
\subsection{Solución}
\[
\|x+y\|^2
=\|x\|^2+2x\cdot y+\|y\|^2.
\]
Por Cauchy--Schwarz,
\[
x\cdot y\leq \|x\|\|y\|.
\]
Así,
\[
\|x+y\|^2
\leq
\|x\|^2+2\|x\|\|y\|+\|y\|^2
=(\|x\|+\|y\|)^2.
\]
Por tanto,
\[
\boxed{\|x+y\|\leq\|x\|+\|y\|}.
\]

\section{Ejercicio 4}
Sea
\[
d(x,y)=|x_1-y_1|+|x_2-y_2|.
\]
\subsection{Solución}
\[
d(x,y)\geq0,
\]
\[
d(x,y)=0\iff x=y,
\]
\[
d(x,y)=d(y,x).
\]
Además,
\[
|x_1-y_1|\leq|x_1-z_1|+|z_1-y_1|,
\]
\[
|x_2-y_2|\leq|x_2-z_2|+|z_2-y_2|.
\]
Sumando,
\[
d(x,y)\leq d(x,z)+d(z,y).
\]
Por tanto,
\[
\boxed{d\text{ es una distancia en }\R^2}.
\]
La bola abierta de centro \((0,0)\) y radio \(1\) es
\[
\boxed{
B((0,0),1)=\{(x,y)\in\R^2:|x|+|y|<1\}.
}
\]

\section{Ejercicio 5}
\subsection{Solución}
\[
\boxed{
B((0,1),2)
=
\{(x,y)\in\R^2:x^2+(y-1)^2<4\}.
}
\]
\[
d((0,1),(0,1))=0<2,
\]
\[
d((0,1),(1,1))=1<2,
\]
\[
d((0,1),(0,3))=2.
\]
Por tanto,
\[
\boxed{
(0,1),(1,1)\in B((0,1),2),
\qquad
(0,3)\notin B((0,1),2).
}
\]

\section{Ejercicio 6}
\subsection{Solución}
Para
\[
A=[1,2]\times[2,3],
\]
\[
\boxed{\Int(A)=(1,2)\times(2,3)}
\]
y
\[
\boxed{\Fr(A)=A\setminus\Int(A)}.
\]
Como \(A\) es cerrado,
\[
\boxed{\Ext(A)=\R^2\setminus A}.
\]
Por tanto,
\[
\boxed{A\text{ es cerrado y no abierto}.}
\]

Para
\[
B=\{(x,y)\in\R^2:4<x^2\leq9\},
\]
\[
B=\left([-3,-2)\cup(2,3]\right)\times\R.
\]
Entonces
\[
\boxed{
\Int(B)=\left((-3,-2)\cup(2,3)\right)\times\R
}
\]
\[
\boxed{
\Fr(B)=\{-3,-2,2,3\}\times\R
}
\]
y
\[
\boxed{
\Ext(B)=
\bigl((-\infty,-3)\cup(-2,2)\cup(3,\infty)\bigr)\times\R.
}
\]
Por tanto,
\[
\boxed{B\text{ no es abierto ni cerrado}.}
\]

\section{Ejercicio 7}
\subsection{Solución}
Para
\[
A=\{(x,y,z):0\leq x,y,z\leq1\},
\]
\[
\boxed{
\Int(A)=\{(x,y,z):0<x<1,\ 0<y<1,\ 0<z<1\}
}
\]
y
\[
\boxed{\Fr(A)=A\setminus\Int(A)}.
\]

Para
\[
B=\{(x,y,z):x^2+y^2\leq2,\ 0\leq z\leq2\},
\]
\[
\boxed{
\Int(B)=\{(x,y,z):x^2+y^2<2,\ 0<z<2\}.
}
\]
Además,
\[
\boxed{
\begin{aligned}
\Fr(B)
={}&\{(x,y,z):x^2+y^2=2,\ 0\leq z\leq2\}\\
&\cup\{(x,y,0):x^2+y^2\leq2\}\\
&\cup\{(x,y,2):x^2+y^2\leq2\}.
\end{aligned}
}
\]

\section{Ejercicio 8}
Sea
\[
A=
\left\{
\left(-\frac1n,\frac1n\right):n\in\mathbb N
\right\}.
\]
\subsection{Solución}
\[
\boxed{\Int(A)=\varnothing}.
\]
Como
\[
\left(-\frac1n,\frac1n\right)\longrightarrow(0,0),
\]
se tiene
\[
\boxed{\Ac(A)=\{(0,0)\}}
\]
\[
\boxed{\Adh(A)=A\cup\{(0,0)\}}
\]
\[
\boxed{\Fr(A)=A\cup\{(0,0)\}}
\]
y
\[
\boxed{
\Ext(A)=\R^2\setminus\left(A\cup\{(0,0)\}\right).
}
\]

\section{Ejercicio 9}
\subsection{Solución}
\[
A_n=
\left(-\frac1n,\frac1n\right)
\times
\left(-\frac1n,\frac1n\right).
\]
Por tanto,
\[
\boxed{
A=\bigcap_{n=1}^{\infty}A_n=\{(0,0)\}.
}
\]
Además, los \(A_n\) son decrecientes, por lo que
\[
\boxed{
B=\bigcup_{n=1}^{\infty}A_n=A_1=(-1,1)\times(-1,1).
}
\]
Así,
\[
\boxed{A\text{ es cerrado y no abierto}}
\]
y
\[
\boxed{B\text{ es abierto y no cerrado}.}
\]

\section{Ejercicio 10}
\subsection{Solución}
Para
\[
A=\{(x,y):x^2+y^2<1\},
\]
\[
\boxed{
A\text{ es acotado},
\qquad
\operatorname{diam}(A)=2.
}
\]
Para
\[
B=\{(x,y,z):x^2+y^2+z^2<1\},
\]
\[
\boxed{
B\text{ es acotado},
\qquad
\operatorname{diam}(B)=2.
}
\]
Para
\[
C=\{(x,y):y=x^2\},
\]
\[
\|(x,x^2)\|\longrightarrow\infty
\qquad(x\to\infty),
\]
por lo que
\[
\boxed{C\text{ no es acotado}.}
\]
Para
\[
D=
\left\{
\left(\frac1n,\frac1n\right):n\in\mathbb N
\right\},
\]
\[
\left\|
\left(\frac1n,\frac1n\right)
\right\|
=
\frac{\sqrt2}{n}\leq\sqrt2.
\]
Por tanto,
\[
\boxed{
D\text{ es acotado},
\qquad
\operatorname{diam}(D)=\sqrt2.
}
\]

\section{Ejercicio 11}
\subsection{Solución}
Para
\[
A=\{(x,y):|x|\leq1,\ |y|\leq1\},
\]
el conjunto es cerrado y acotado. Por tanto,
\[
\boxed{A\text{ es compacto}.}
\]
Para
\[
B=\{(x,y):xy\leq1\},
\]
se tiene
\[
(0,n)\in B\qquad\forall n\in\mathbb N,
\]
por lo que no es acotado. Así,
\[
\boxed{
B\text{ no es compacto y no existe un compacto que lo contenga}.
}
\]
Para
\[
C=\{(x,y,z):0<x^2+y^2+z^2\leq1\},
\]
el conjunto es acotado pero no cerrado, ya que \((0,0,0)\) es punto de acumulación y no pertenece a \(C\). Por tanto,
\[
\boxed{C\text{ no es compacto}.}
\]
Su clausura es
\[
\boxed{
\overline C=
\{(x,y,z):x^2+y^2+z^2\leq1\},
}
\]
que es el compacto mínimo que contiene a \(C\).

\section{Ejercicio 12}
\subsection{Solución}
Si
\[
A=(2,-1),
\qquad
v=(7,6),
\]
el punto final \(B\) cumple
\[
\overrightarrow{AB}=B-A=v.
\]
Por tanto,
\[
B=A+v=(9,5).
\]
\[
\boxed{B=(9,5)}
\]

\section{Ejercicio 13}
\subsection{Solución}
El vector buscado es
\[
w=-\frac12(-2,4)=(1,-2).
\]
Si \(P\) es el punto inicial y \(Q=(2,0)\) el final,
\[
Q-P=w.
\]
Entonces
\[
P=Q-w=(2,0)-(1,-2)=(1,2).
\]
\[
\boxed{P=(1,2)}
\]

\section{Ejercicio 14}
\subsection{Solución}
\[
\overrightarrow{AB}
=
(-3,2)-(2,1)
=
(-5,1).
\]
Por tanto,
\[
\|\overrightarrow{AB}\|
=
\sqrt{25+1}
=
\sqrt{26}.
\]
\[
\boxed{\|\overrightarrow{AB}\|=\sqrt{26}}
\]

\section{Ejercicio 15}
\subsection{Solución}
Para
\[
a=(4,-3),
\qquad
\|a\|=5,
\]
\[
\boxed{
u_a=
\left(\frac45,-\frac35\right).
}
\]
Para
\[
b=(1,-4,2),
\qquad
\|b\|=\sqrt{21},
\]
\[
\boxed{
u_b=
\frac1{\sqrt{21}}(1,-4,2).
}
\]

\section{Ejercicio 16}
\subsection{Solución}
\[
a\times b
=
\begin{vmatrix}
i&j&k\\
1&2&3\\
3&2&1
\end{vmatrix}.
\]
Por tanto,
\[
\boxed{
a\times b=(-4,8,-4).
}
\]

\section{Ejercicio 17}
\subsection{Solución}
\[
a\times b
=
\begin{vmatrix}
i&j&k\\
2&0&1\\
1&-1&3
\end{vmatrix}
=
(1,-5,-2).
\]
Así,
\[
\boxed{a\times b=(1,-5,-2)}
\]
y
\[
\boxed{b\times a=(-1,5,2)}.
\]

\section{Ejercicio 18}
\subsection{Solución}
Sean
\[
x=(x_1,\ldots,x_n),
\quad
y=(y_1,\ldots,y_n),
\quad
z=(z_1,\ldots,z_n).
\]
Entonces
\[
x\cdot y
=
\sum_{i=1}^nx_iy_i
=
\sum_{i=1}^ny_ix_i
=
y\cdot x.
\]
\[
\boxed{x\cdot y=y\cdot x}
\]
Además,
\[
x\cdot(y+z)
=
\sum_{i=1}^nx_i(y_i+z_i)
=
x\cdot y+x\cdot z.
\]
\[
\boxed{x\cdot(y+z)=x\cdot y+x\cdot z}
\]
Y
\[
x\cdot(\lambda y)
=
\sum_{i=1}^nx_i\lambda y_i
=
\lambda(x\cdot y).
\]
Por conmutatividad,
\[
y\cdot(\lambda x)=\lambda(x\cdot y).
\]
Así,
\[
\boxed{
x\cdot(\lambda y)
=
y\cdot(\lambda x)
=
\lambda(x\cdot y).
}
\]

\section{Ejercicio 19}
\subsection{Solución}
\[
a\times b
=
\begin{vmatrix}
i&j&k\\
a_1&a_2&a_3\\
b_1&b_2&b_3
\end{vmatrix}.
\]
Intercambiando las dos últimas filas,
\[
\boxed{a\times b=-(b\times a)}.
\]
Por linealidad del determinante respecto de una fila,
\[
a\times(b+c)
=
a\times b+a\times c.
\]
Por tanto,
\[
\boxed{
a\times(b+c)=a\times b+a\times c.
}
\]
Asimismo,
\[
a\times(\lambda b)=\lambda(a\times b)
\]
y
\[
(\lambda a)\times b=\lambda(a\times b).
\]
Así,
\[
\boxed{
a\times(\lambda b)
=
(\lambda a)\times b
=
\lambda(a\times b).
}
\]

\section{Ejercicio 20}
\subsection{Solución}
Tomamos
\[
a=(1,0,0),
\qquad
b=(0,1,0),
\qquad
c=(0,1,0).
\]
Entonces
\[
a\times b=(0,0,1)
\]
y
\[
(a\times b)\times c
=
(0,0,1)\times(0,1,0)
=
(-1,0,0).
\]
Sin embargo,
\[
b\times c=0,
\]
por lo que
\[
a\times(b\times c)=0.
\]
Por tanto,
\[
\boxed{
(a\times b)\times c
\neq
a\times(b\times c).
}
\]

\section{Ejercicio 21}
\subsection{Solución}
\[
r=-6\cos\theta.
\]
Multiplicando por \(r\),
\[
r^2=-6r\cos\theta.
\]
Como
\[
r^2=x^2+y^2,
\qquad
r\cos\theta=x,
\]
se obtiene
\[
\boxed{x^2+y^2+6x=0}
\]
o
\[
\boxed{(x+3)^2+y^2=9}.
\]
Una parametrización es
\[
\boxed{
\begin{cases}
x=-6\cos^2\theta,\\
y=-6\sen\theta\cos\theta,
\end{cases}
\qquad
0\leq\theta<2\pi.
}
\]

\section{Ejercicio 22}
\subsection{Solución}
La curva pasa por el origen cuando
\[
r=0.
\]
Entonces
\[
1-2\sen\theta=0
\]
y
\[
\sen\theta=\frac12.
\]
Por tanto,
\[
\boxed{
\theta=\frac{\pi}{6},
\qquad
\theta=\frac{5\pi}{6}.
}
\]
La curva \(r=1-2\sen\theta\) es un limacón con lazo interior, simétrico respecto del eje \(Y\).

\section{Ejercicio 23}
\subsection{Solución}
La curva
\[
\boxed{r=2\sen(3\theta)}
\]
es una rosa polar de tres pétalos de longitud \(2\).

La curva
\[
\boxed{r^2=\cos(2\theta)}
\]
es una lemniscata de Bernoulli con los dos lóbulos sobre el eje \(X\).

\section{Ejercicio 24}
\subsection{Solución}
Para los puntos distintos del origen,
\[
\cos\theta=\cos(2\theta).
\]
Como
\[
\cos(2\theta)=2\cos^2\theta-1,
\]
\[
2\cos^2\theta-\cos\theta-1=0.
\]
Factorizando,
\[
(2\cos\theta+1)(\cos\theta-1)=0.
\]
Así,
\[
\cos\theta=1
\]
o
\[
\cos\theta=-\frac12.
\]
Se obtienen los puntos
\[
(1,0),
\]
\[
\left(\frac14,\frac{\sqrt3}{4}\right),
\qquad
\left(\frac14,-\frac{\sqrt3}{4}\right).
\]
Además, ambas curvas pasan por el origen. Por tanto,
\[
\boxed{
(0,0),\quad
(1,0),\quad
\left(\frac14,\frac{\sqrt3}{4}\right),\quad
\left(\frac14,-\frac{\sqrt3}{4}\right).
}
\]

\section{Ejercicio 25}
Sea
\[
f(x,y)=\frac{x+y}{|x+y|}.
\]
\subsection{Solución}
Debe cumplirse
\[
x+y\neq0.
\]
Por tanto,
\[
\boxed{
\Dom(f)=\{(x,y)\in\R^2:x+y\neq0\}.
}
\]
Además,
\[
f(x,y)=
\begin{cases}
1,&x+y>0,\\
-1,&x+y<0.
\end{cases}
\]
Por tanto,
\[
\boxed{\Ima(f)=\{-1,1\}}.
\]

\section{Ejercicio 26}
Sean
\[
f(x,y)=\frac{\sqrt{4-y^2}}{3+\sqrt{2-x^2}}
\]
y
\[
g(x,y,z)=\frac{3z}{y^2-x^2}.
\]
\subsection{Solución}
Para \(f\),
\[
4-y^2\geq0,
\qquad
2-x^2\geq0.
\]
Por tanto,
\[
\boxed{
\Dom(f)
=
[-\sqrt2,\sqrt2]\times[-2,2].
}
\]
Como
\[
0\leq\sqrt{4-y^2}\leq2
\]
y
\[
3\leq3+\sqrt{2-x^2}\leq3+\sqrt2,
\]
se obtiene
\[
\boxed{
\Ima(f)=\left[0,\frac23\right].
}
\]

Para \(g\),
\[
y^2-x^2\neq0.
\]
Así,
\[
\boxed{
\Dom(g)
=
\{(x,y,z)\in\R^3:y\neq x,\ y\neq-x\}.
}
\]
Para cualquier \(k\in\R\), tomando
\[
x=0,\qquad y=1,\qquad z=\frac{k}{3},
\]
se tiene \(g(0,1,k/3)=k\). Por tanto,
\[
\boxed{\Ima(g)=\R}.
\]

\section{Ejercicio 27}
Sea
\[
f(x,y)
=
xy\sen\left(
\frac1{\sqrt{x^2+y^2}}
\right).
\]
\subsection{Solución}
Debe cumplirse
\[
x^2+y^2>0.
\]
Por tanto,
\[
\boxed{
\Dom(f)=\R^2\setminus\{(0,0)\}.
}
\]

\section{Ejercicio 28}
Sea
\[
f(x,y)
=
\left(
x^2+y^2,\,
\sen(xy),\,
\ln(x+y)
\right).
\]
\subsection{Solución}
La única restricción procede del logaritmo:
\[
x+y>0.
\]
Por tanto,
\[
\boxed{
\Dom(f)=\{(x,y)\in\R^2:x+y>0\}.
}
\]

\section{Ejercicio 29}
\subsection{Solución}
Para
\[
f(x,y)=x^2+y^2,
\]
las curvas de nivel son
\[
x^2+y^2=k.
\]
Por tanto,
\[
\boxed{
C_k=
\begin{cases}
\varnothing,&k<0,\\
\{(0,0)\},&k=0,\\
\text{circunferencia de radio }\sqrt{k},&k>0.
\end{cases}
}
\]

Para
\[
g(x,y)=x+y,
\]
las curvas de nivel son
\[
\boxed{
C_k:\ y=-x+k,
\qquad k\in\R.
}
\]

\section{Ejercicio 30}
\subsection{Solución}
Las curvas de nivel satisfacen
\[
xy=k.
\]
Si \(k\neq0\),
\[
y=\frac{k}{x},
\]
que son hipérbolas equiláteras.

Si \(k=0\),
\[
xy=0,
\]
por lo que
\[
x=0
\]
o
\[
y=0.
\]
Así,
\[
\boxed{
C_k=
\begin{cases}
\left\{(x,y):y=\dfrac{k}{x}\right\},&k\neq0,\\[2mm]
\{x=0\}\cup\{y=0\},&k=0.
\end{cases}
}
\]

\end{document}
